% =====================================================================
%  CHAPTER 2 — BACKGROUND AND RELATED WORK
%  Graduation thesis: A hybrid (signature + AI) web intrusion detection
%  system operating on Apache HTTP/HTTPS access logs.
%
%  HOW TO USE
%  ----------
%  * Standalone build: this file compiles on its own (pdflatex twice to
%    resolve references and the embedded bibliography).
%  * Integration into your thesis: delete the lines between the two
%    "STANDALONE WRAPPER" markers (the \documentclass...\begin{document}
%    block and the matching \end{document}), keep everything from
%    "\chapter{...}" to the end of the thebibliography environment, and
%    move the \bibitem entries into your global bibliography (.bib).
% =====================================================================

% --------------------- STANDALONE WRAPPER (begin) --------------------
\documentclass[12pt,a4paper]{report}
\usepackage[utf8]{inputenc}
\usepackage[T1]{fontenc}
\usepackage{lmodern}
\usepackage{amsmath}
\usepackage{booktabs}
\usepackage{array}
\usepackage{enumitem}
\usepackage{graphicx}
\usepackage{url}
\usepackage[hidelinks]{hyperref}
\usepackage{geometry}
\geometry{margin=1in}
\usepackage{textcomp}

% Roman-numeral enumerated lists used for the "scientific" listings.
\newlist{romanlist}{enumerate}{1}
\setlist[romanlist]{label=(\roman*), leftmargin=2.2em, itemsep=2pt}

\begin{document}
% ---------------------- STANDALONE WRAPPER (end) ---------------------

\chapter{Background and Related Work}
\label{chap:background}

This chapter establishes the context, motivation, and the core knowledge
that the thesis relies upon. Section~\ref{sec:scope} first delimits the
scope of the web intrusion detection problem that the system targets: its
input data, its expected output, and the classes of attacks that fall
inside and outside the processing boundary. Section~\ref{sec:related}
then surveys representative studies that address attack detection on the
same HTTP dataset, analysing their strengths and weaknesses in order to
expose the research gap and the motivation behind this work. The
remaining sections present the foundational knowledge most tightly
connected to the system: the notion of an intrusion detection system and
its two detection paradigms (Section~\ref{sec:ids}), common web
application attacks (Section~\ref{sec:webattacks}), signature-based
detection with regular expressions (Section~\ref{sec:signature}),
supervised machine learning for classification (Section~\ref{sec:sup}),
unsupervised anomaly detection following the one-class learning principle
(Section~\ref{sec:unsup}), and finally the feature extraction technique
together with the evaluation metrics used throughout the thesis
(Section~\ref{sec:features}). These topics are the direct basis for the
three-tier hybrid architecture realised in the subsequent chapters.

% =====================================================================
\section{Scope of Research}
\label{sec:scope}

This thesis focuses on the problem of \emph{web intrusion detection} at
the application layer through the analysis of access logs produced by a
web server. The task is framed as a binary classification problem: for
every recorded HTTP/HTTPS request, the system must decide whether the
request is \emph{benign} (clean) or an \emph{attack}, and it must
additionally rank the severity of each alert so that a security analyst
can prioritise the investigation.

Regarding the input data, the system operates directly on the widely used
Apache log formats, namely the Combined Log Format and the Common Log
Format. Each log line carries fields such as the source IP address, the
timestamp, the request line (method, URL, and protocol version), the
returned status code, the \texttt{Referer} header, and the
\texttt{User-Agent} header. During implementation, the log parsing module
(\texttt{parse\_log\_entry} in \texttt{apache\_log.py}) decomposes the
request line into method, path, query string, and parameters, and it
iteratively URL-decodes the request to recover the original payload before
feeding it to the detection tiers. A noteworthy characteristic is that the
body of \texttt{POST} requests is also preserved and appended to the URL
as a \texttt{POST\_BODY} parameter, so that parameter-analysis techniques
apply to attacks delivered through the request body and not merely through
the query string.

Regarding the output, for every log line the system produces a structured
record (in JSON Lines format) containing the \texttt{is\_attack} verdict,
an aggregated threat score, a severity level (\textsc{critical},
\textsc{high}, \textsc{medium}, \textsc{low}), and the list of detection
tiers that fired. The alerts are aggregated into a visual dashboard and a
live alert file for real-time monitoring.

The classes of attacks addressed by the thesis are application-layer web
attacks that manifest directly in the content of the HTTP request.
Specifically, the system targets the following attack groups:

\begin{romanlist}
  \item SQL Injection;
  \item Cross-Site Scripting (XSS);
  \item Path Traversal and Local File Inclusion;
  \item Operating-system Command Injection;
  \item Remote File Inclusion;
  \item Access to sensitive configuration or backup files;
  \item Use of automated scanning tools identified through the
        \texttt{User-Agent} header;
  \item Parameter Tampering and HTTP Method Tampering.
\end{romanlist}

Conversely, several topics lie outside the scope of this thesis. The
system does not handle network- or transport-layer attacks such as
volumetric distributed denial of service, it does not perform packet-level
decryption of TLS traffic (it works on logs already decrypted by the
server), and it does not assume the blocking role of a Web Application
Firewall. The emphasis of the thesis is on \emph{detection} and
\emph{alerting}, which form the prerequisite for subsequent response
measures.

For training and evaluation, the thesis builds a balanced dataset of
$111{,}065$ log lines, of which one half is attack traffic and the other
half is benign traffic. This dataset combines three sources: the standard
HTTP DATASET CSIC 2010~\cite{csic2010}, real attack samples extracted from
the open-source PayloadsAllTheThings repository, and synthetic traffic
generated automatically to increase the diversity of encoding variants.
Choosing CSIC 2010 as the backbone of the dataset allows the results of
this thesis to be compared directly with the related studies presented in
the next section, since that dataset is also the one used by most of those
studies.

% =====================================================================
\section{Related Work}
\label{sec:related}

Enhancing intrusion detection with machine learning has been studied
extensively. This section reviews representative studies that operate on
the HTTP CSIC 2010 dataset, analyses the strengths and weaknesses of each
approach, and thereby identifies the gap that this thesis aims to fill.

\subsection{Comparing machine learning algorithms on CSIC 2010
(Pham et al., 2016)}

Pham, Hoang, and Vu~\cite{pham2016} surveyed and compared five machine
learning algorithms---Random Forest, Logistic Regression, Decision Tree,
AdaBoost, and Stochastic Gradient Descent---on the CSIC 2010 dataset. The
strength of their approach lies in the feature extraction step applied to
the raw text of the HTTP request: the authors combine three techniques,
namely tokenizing, occurrence counting, and re-weighting by the inverse
document frequency (TF--IDF), producing a sparse matrix as input to the
classifier. The experimental results show that Logistic Regression
achieved the best performance, with an F1 score of about $96.11\%$ for the
attack class; after tuning the inverse regularisation strength (choosing
$C = 13.0$) together with recursive feature elimination (retaining
$77{,}558$ features), the precision and recall were further improved.

Nevertheless, this approach exhibits several limitations. First, the
bag-of-words representation with tens of thousands of features depends
heavily on the vocabulary learned from the training set, which makes the
model bulky and prone to losing validity when faced with traffic from a
different web application that uses a different vocabulary. Second, this is
a purely \emph{single-model supervised} approach that requires fully
labelled data and, by nature, only learns the characteristics of attacks
already seen; the authors themselves acknowledge in their conclusion that
evaluation on unseen attacks remains open. Third, the study is conducted
in batch mode on a pre-processed dataset and does not address operation
directly on server logs in real time.

\subsection{Analysing HTTP requests with structural features
(Althubiti et al., 2017)}

Althubiti, Yuan, and Esterline~\cite{althubiti2017} also classify HTTP
requests on CSIC 2010, but under a different feature philosophy. Instead of
a text representation, the authors use nine \emph{structural}, numeric
features proposed by Nguyen et al.~\cite{nguyen2011}, such as the length of
the request, the length of the argument section, the number of arguments,
the length of the path, and the number of special characters in the path.
They employ the Weka toolkit with the \texttt{WrapperSubsetEval} attribute
evaluator (using the J48 classifier and ten-fold cross-validation) combined
with the \texttt{BestFirst} search strategy to select the five best
features. On this compact feature set, most algorithms (Random Forest,
Logistic Regression, J48, AdaBoost, SGD) achieve a very high detection
rate of approximately $99.9\%$, with only Na\"{\i}ve Bayes performing
markedly worse (around $88\%$).

The advantage of this work is a lightweight numeric feature set that is
easy to compute, fast to train, and highly accurate; feature selection
further reduces the training time while improving the results. However,
the near-perfect detection figures also raise questions about
generalisation when the training and test data share the same
distribution, and the study remains entirely within the supervised
paradigm with no mechanism for novel attacks. Moreover, although features
based purely on length and counts are sensitive to buffer-overflow
attacks, they may miss attacks that are semantically sophisticated yet not
anomalous in length.

\subsection{Comparing seven algorithms and a voting classifier
(Rawashdeh, 2023)}

Rawashdeh~\cite{rawashdeh2023} extends the comparison to seven
algorithms---Decision Tree, Random Forest, Gradient Boosting, XGBoost,
AdaBoost, Multilayer Perceptron, and a Voting Classifier---on the Kaggle
edition of CSIC 2010 ($61{,}065$ samples, $16$ features corresponding to
HTTP header fields). The non-numeric features are converted to numeric
form through Label Encoding. The results show that the Voting Classifier
and the Decision Tree deliver the best performance, with an accuracy of
about $89\%$. A notable contribution of this work is its demonstration of
the benefit of \emph{ensemble} learning through a voting mechanism, in
which the final label is decided by the majority vote of the component
classifiers.

The limitation of the study is that its overall accuracy ($\approx 89\%$)
is lower than that of the two works above, partly because label-encoding
categorical attributes does not preserve the semantics of the HTTP
content. Furthermore, although a voting mechanism is used, all components
are supervised models, so the system still inherits the common drawback of
being unable to detect attack types lying outside the training
distribution.

\subsection{Deep learning and one-class approaches on the same dataset}

Beyond classical classifiers, several studies surveyed
in~\cite{rawashdeh2023} apply deep learning to CSIC 2010. Kim et
al.~\cite{kim2020} combine a Convolutional Neural Network with a Long
Short-Term Memory network (CNN with LSTM) and reach an accuracy of about
$91.54\%$. Of particular relevance to this thesis, Vartouni et
al.~\cite{vartouni2018} use a Stacked Auto-Encoder---a reconstruction-based
model with an \emph{anomaly detection} flavour---and obtain an accuracy of
about $88.32\%$. Betarte et al.~\cite{betarte2018} compare Random Forest,
K-Nearest Neighbours, and Support Vector Machine, with Random Forest
performing best in that group. These works reveal two trends: on the one
hand, complex deep models do not necessarily outperform compact structural
classifiers on CSIC 2010; on the other hand, the reconstruction-based
anomaly direction opens up the possibility of detecting novel attacks
without requiring attack labels during training.

\subsection{Synthesis and research motivation}

Table~\ref{tab:related} summarises the representative studies just
reviewed, contrasting them along the criteria of dataset, method, approach,
and main limitation. Several important commonalities emerge from this
table.

\begin{table}[ht]
  \centering
  \caption{Summary of related work on web attack detection over CSIC 2010.}
  \label{tab:related}
  \renewcommand{\arraystretch}{1.45}
  \setlength{\tabcolsep}{8pt}
  \footnotesize
  \begin{tabular}{|p{2.3cm}|p{1.9cm}|p{3.0cm}|p{1.9cm}|p{2.8cm}|}
    \hline
    \textbf{Study} & \textbf{Dataset} & \textbf{Representative method} &
    \textbf{Approach} & \textbf{Main limitation} \\
    \hline
    Pham et al., 2016 & CSIC 2010 &
    LR, RF, DT, AdaBoost, SGD on TF--IDF features &
    Supervised, single model &
    Bulky features, no novel-attack handling \\
    \hline
    Nguyen et al., 2011 & CSIC 2010, ECML/PKDD &
    GeFS feature selection (CFS, mRMR) & Supervised &
    Hand-crafted structural features \\
    \hline
    Althubiti et al., 2017 & CSIC 2010 &
    RF, LR, J48, AdaBoost, SGD, NB on 5 features &
    Supervised & Near-perfect scores, generalisation doubts \\
    \hline
    Rawashdeh, 2023 & CSIC 2010 (Kaggle) &
    DT, RF, GB, XGB, AdaBoost, MLP, Voting & Supervised ensemble &
    Accuracy $\sim$89\%, label encoding loses semantics \\
    \hline
    Kim et al., 2020 & CSIC 2010 & CNN + LSTM, DNN &
    Supervised deep learning & Complex, needs much labelled data \\
    \hline
    Vartouni et al., 2018 & CSIC 2010 &
    Stacked Auto-Encoder & Anomaly detection &
    Accuracy $\sim$88\%, recall threshold-dependent \\
    \hline
  \end{tabular}
\end{table}

First, the vast majority of studies adopt a \emph{single} approach,
either purely supervised learning or purely anomaly detection. The
supervised direction yields high accuracy on known attacks but requires
labelled data and is weak against entirely new attacks. The anomaly
direction has the potential to catch previously unseen attacks but
typically suffers from a high false-alarm rate and poor stability across
algorithms. Second, most studies work on offline pre-processed feature
sets; few systems operate directly on raw server logs and handle practical
situations such as multi-layer encoded payloads or attacks delivered
through the \texttt{POST} body. Third, the signature-based detection
layer---which has very high precision and good explainability---is often
ignored or kept separate from the machine learning models rather than
combined in a systematic way.

These observations are precisely the motivation of this thesis. The
system's own experiments confirm the complementarity of the approaches:
the signature tier reaches a precision of $99.61\%$ but a recall of only
$24.07\%$ (missing most attacks), whereas the unsupervised anomaly models
attain high recall but widely varying precision across algorithms.
Therefore, the thesis proposes a \emph{three-tier hybrid} architecture
that simultaneously combines (i) signature- and rule-based detection,
(ii) supervised learning with Random Forest, and (iii) unsupervised
anomaly detection following the one-class learning principle. A smart
consensus decision mechanism lets the supervised tier validate the
detections of the anomaly tier, thereby exploiting the high recall of
anomaly detection while keeping false alarms under control. This is the
core difference between this thesis and the surveyed studies.

% =====================================================================
\section{Intrusion Detection Systems}
\label{sec:ids}

An Intrusion Detection System (IDS) is a security component whose task is
to monitor the events occurring in a computer system or network and to
analyse them for signs of intrusive behaviour~\cite{tsai2009}. According to
the monitoring scope, an IDS is usually divided into a Network-based IDS,
which monitors traffic crossing the network, and a Host-based IDS, which
monitors activity on a specific machine. This thesis belongs to the web
application intrusion detection category, a specialised form that operates
on the application-layer logs of a web server.

According to the detection method, the classical literature distinguishes
two main paradigms, summarised by Althubiti et al.~\cite{althubiti2017} as
follows. The first paradigm is \emph{misuse detection}, also called
signature-based detection, which operates by matching the current activity
against known attack patterns, usually through pattern-matching
algorithms. Its advantage is high accuracy and few false alarms for known
attacks, but its inherent drawback is the inability to detect new attacks
that are absent from the signature set. The second paradigm is
\emph{anomaly detection}, which studies the behaviour of the user or the
system in order to build a model of ``normal'' behaviour, then treats
significant deviations from this model as anomalies; this direction is
typically based on machine learning techniques and is able to detect new
attacks, at the cost of potentially generating more false alarms.

These two paradigms are clearly complementary: misuse detection is strong
in precision while anomaly detection is strong in covering unknown
situations. This observation is the theoretical foundation for the hybrid
architecture of the thesis, in which the signature tier
(Section~\ref{sec:signature}) represents misuse detection and the
unsupervised tier (Section~\ref{sec:unsup}) represents anomaly detection;
the supervised tier (Section~\ref{sec:sup}) plays an intermediary role,
both generalising from known attacks and providing a confidence
probability that reconciles the two paradigms.

% =====================================================================
\section{Common Web Application Attacks}
\label{sec:webattacks}

To design appropriate detection rules and features, it is necessary to
understand the nature of the web attacks that the system targets. Most of
these attacks appear in the list of the ten most critical web application
security risks published by the OWASP project~\cite{owasp2021}. In the
OWASP Top 10 of 2021, the ``Injection'' group ranks at position A03 and
includes several sub-types such as SQL injection, operating-system command
injection, and cross-site scripting~\cite{owasp2021}. The attack types
directly relevant to the thesis are presented below.

\begin{romanlist}
  \item \textbf{SQL Injection.} The attack occurs when untrusted data is
        sent to an interpreter as part of a database query~\cite{pham2016}.
        By inserting clauses such as \texttt{UNION SELECT}, \texttt{OR
        '1'='1}, or statement terminators, the attacker can trick the
        interpreter into executing unintended operations and thus access or
        modify data without authorisation.

  \item \textbf{Cross-Site Scripting (XSS).} An XSS flaw appears when an
        application takes untrusted data and sends it to a browser without
        proper validation or escaping~\cite{pham2016}. The attacker can
        then execute script code in the victim's browser to hijack the
        session, deface the page, or redirect the user to a malicious site.

  \item \textbf{Path Traversal and Local File Inclusion.} The attacker
        exploits the \texttt{../} sequence to escape the intended directory
        and access sensitive system files such as \texttt{/etc/passwd} or
        \texttt{win.ini}.

  \item \textbf{Command Injection.} The attack allows the execution of
        arbitrary operating-system commands through command-chaining
        characters (\texttt{;}, \texttt{|}, \texttt{\&\&}, the backtick)
        placed before dangerous commands such as \texttt{rm}, \texttt{cat},
        \texttt{wget}, or \texttt{bash}.

  \item \textbf{Remote File Inclusion.} An attack in which a parameter
        value is a remote URL (starting with \texttt{http://},
        \texttt{https://}, or \texttt{ftp://}), aiming to force the server
        to load and execute code from an external source.

  \item \textbf{Sensitive file access, parameter tampering, and method
        tampering.} Sensitive file access is an attempt to read
        configuration or backup files such as \texttt{.env},
        \texttt{wp-config.php}, or \texttt{.sql}. Parameter tampering is the
        modification of a parameter name relative to the legitimate names an
        endpoint usually accepts, for example changing \texttt{modo} to
        \texttt{modoA}. HTTP Method Tampering is the use of unusual methods
        such as \texttt{PUT}, \texttt{DELETE}, or \texttt{TRACE} on
        endpoints that normally serve only \texttt{GET}/\texttt{POST}.
\end{romanlist}

Enumerating these attack classes helps the reader grasp the scope of the
problem; how each one is actually detected---through dedicated regular
expressions and rules---is detailed in the next chapter.

% =====================================================================
\section{Signature-Based Detection with Regular Expressions}
\label{sec:signature}

Signature-based detection is the technique of matching the content of a
request against a set of known attack patterns. In the thesis, these
patterns are expressed as \emph{regular expressions} (regex)---a formal
language that allows a compact description of sets of character strings.
Each attack group in Section~\ref{sec:webattacks} corresponds to a
pre-compiled regular expression, for example \texttt{SQLI\_RE} for SQL
injection or \texttt{XSS\_RE} for cross-site scripting. When a request is
loaded, the \texttt{detect\_rule\_based} function applies these expressions
to the payload in turn and returns the list of rules that fired.

An important challenge for signature-based detection is the \emph{encoding
evasion} technique, in which attackers encode the URL in several layers to
hide the payload, so that matching on the raw string fails. To overcome
this, the system iteratively URL-decodes the request up to three times (the
\texttt{multi\_decode} function and the decoding loop in
\texttt{detect\_rule\_based}) to recover the original payload---for
instance the tautology \texttt{'OR'a'='a}---before matching it against the
rules.

The advantage of the signature tier is its very high precision and its
transparent explainability: every alert explicitly indicates which rule
fired, helping the analyst quickly verify it. However, as noted, the
inherent weakness of this direction is low recall---it catches only what
has been explicitly described in the rule set. The thesis's experimental
results reflect exactly this property, with a precision of $99.61\%$ but a
recall of only $24.07\%$. This is precisely why the signature tier must be
complemented by the machine learning tiers in the following sections.

% =====================================================================
\section{Supervised Machine Learning for Attack Classification}
\label{sec:sup}

Supervised learning is an approach in which the model learns from a
labelled dataset, each sample consisting of a feature vector and a class
label (here, \emph{attack} or \emph{benign}), in order to build a function
that maps features to labels and predict for new data~\cite{pham2016}. The
thesis uses supervised learning as the second detection tier, with two
algorithms implemented in \texttt{supervised\_detector.py}: Random Forest
(serving as the primary model) and Logistic Regression (serving as a
comparison baseline).

\subsection{Random Forest}

The Random Forest, proposed by Breiman in 2001~\cite{breiman2001}, is an
ensemble learning method that adds randomness to the bagging technique.
Each decision tree in the forest is built from a different bootstrap sample
of the data; moreover, at each node, instead of selecting the best split
over all attributes, the algorithm considers only a randomly chosen subset
of attributes~\cite{pham2016}. This strategy, although seemingly
counter-intuitive, helps the random forest achieve high performance and
robustness against overfitting. For a classification task, the final result
is the class voted for by the most trees.

Formally, when growing each tree, a node is split so as to minimise the
impurity of the resulting child nodes. Writing $p_i$ for the proportion of
class $i$ among the samples in a node, the two impurity measures commonly
used are the Gini impurity and the Shannon entropy,
\begin{equation}
  G = 1 - \sum_{i} p_i^{2},
  \qquad
  H = -\sum_{i} p_i \log_2 p_i .
  \label{eq:gini}
\end{equation}
For an ensemble of $B$ trees $\{T_b\}_{b=1}^{B}$, the predicted class of a
sample $x$ is the majority vote $\hat{y}(x) = \operatorname{mode}\{T_b(x)\}$,
and the probability that $x$ is an attack is estimated as the fraction of
trees voting for the attack class,
\begin{equation}
  p_{\text{attack}}(x) = \frac{1}{B}\sum_{b=1}^{B}
  \mathbf{1}\!\left[T_b(x) = \text{attack}\right],
  \label{eq:rfvote}
\end{equation}
where $\mathbf{1}[\cdot]$ is the indicator function. This probability, rather
than the hard label, is what the smart consensus mechanism later exploits.

In the thesis, the Random Forest is configured with $200$ trees and
balanced class weights (\texttt{class\_weight="balanced"}); being
threshold-based, it requires no feature scaling. Crucially, besides the
predicted label it also returns the \emph{probability} of the attack class
through \texttt{predict\_proba}. This probability is pivotal for the smart
consensus mechanism, letting the supervised tier validate the unsupervised
tier's detections at different confidence levels rather than issuing a hard
decision. In experiments the Random Forest attains an F1 score of about
$94.78\%$, the best single model in the system.

\subsection{Logistic Regression}

Logistic Regression is a linear classifier that estimates the probability
of a sample belonging to a class as a smooth function of the input
features~\cite{pham2016}, trained by maximum likelihood. Concretely, for a
weight vector $\mathbf{w}$ and bias $b$ it models the attack probability with
the logistic (sigmoid) function,
\begin{equation}
  P(y=1 \mid x) = \sigma(\mathbf{w}^{\top}x + b),
  \qquad
  \sigma(z) = \frac{1}{1 + e^{-z}},
  \label{eq:sigmoid}
\end{equation}
and its parameters are fitted by minimising the binary cross-entropy
(log-loss) over the $n$ training samples,
\begin{equation}
  \mathcal{L} = -\sum_{i=1}^{n}
  \Big[\, y_i \log p_i + (1-y_i)\log(1-p_i) \,\Big],
  \qquad p_i = P(y=1 \mid x_i).
  \label{eq:logloss}
\end{equation}
Pham et al.~\cite{pham2016} note that it suits binary classification well and
has low variance, so it is less prone to overfitting. In the thesis it serves
as a backup and reference model; being scale-sensitive, it is applied with
a feature-scaling step (\texttt{StandardScaler}), unlike the Random Forest.
Its F1 score is about $79.99\%$, lower than the Random Forest, consistent
with a linear model's difficulty in capturing the non-linear interactions
among the $22$ features.

An important commonality of both supervised models in the thesis is that
they use the \emph{same 22-feature space} as the unsupervised tier
(Section~\ref{sec:features}). Consequently, the performance difference
between the two tiers reflects the genuine ``labelled versus unlabelled''
distinction rather than being confounded by differences in feature
representation.

% =====================================================================
\section{Unsupervised Anomaly Detection via One-Class Learning}
\label{sec:unsup}

The third detection tier of the thesis follows the unsupervised direction
in order to catch new attacks that the supervised tier has never learned.
The core principle here is \emph{one-class learning}: the model is trained
only on ``normal'' data to learn the distribution of clean traffic, then
treats any significant deviation from that distribution as an anomaly, that
is, an attack~\cite{scholkopf1999}. This approach requires no attack labels
during training, which fits the reality that attack traffic is both rare
and difficult to collect comprehensively.

An important methodological point that the thesis strictly observes is to
\emph{train the third tier only on clean data} ($34{,}941$ benign log
lines), rather than on a mixture of attack and benign traffic. Experiments
show that with incorrect training (on a mixture containing $50\%$ attacks),
the F1 score of the models reaches only about $33\%$; with correct training
following the one-class principle on clean data, the F1 score rises sharply,
for example the Local Outlier Factor model attains $91.57\%$. This
improvement confirms that one-class learning is not merely a technical
choice but a precondition for the unsupervised tier to function correctly.
The thesis implements three complementary algorithms, described below.

\subsection{Local Outlier Factor (LOF)}

The Local Outlier Factor, proposed by Breunig et al. in
2000~\cite{breunig2000}, is an outlier-detection algorithm based on
\emph{local density}. For each data point, LOF estimates the local density
from the distances to its $k$ nearest neighbours, then compares the density
of that point with the density of its neighbours. A point lying in a region
markedly sparser than its neighbourhood receives a high LOF score and is
treated as an outlier. The ``local'' nature of the algorithm enables it to
detect subtle anomalies sitting between regions of different density, which
methods based on a global threshold tend to miss.

Formally, let $N_k(p)$ be the set of the $k$ nearest neighbours of a point
$p$ and let $d(p,o)$ be the distance between $p$ and $o$. The
\emph{reachability distance} of $p$ from $o$ smooths short distances using
the $k$-distance of $o$,
\begin{equation}
  \text{reach-dist}_k(p,o) = \max\{\, k\text{-distance}(o),\; d(p,o) \,\}.
  \label{eq:reachdist}
\end{equation}
The \emph{local reachability density} of $p$ is the inverse of the average
reachability distance to its neighbours,
\begin{equation}
  \text{lrd}_k(p) = \left( \frac{1}{|N_k(p)|}
  \sum_{o \in N_k(p)} \text{reach-dist}_k(p,o) \right)^{-1},
  \label{eq:lrd}
\end{equation}
and the Local Outlier Factor of $p$ is the average ratio of the neighbours'
densities to its own,
\begin{equation}
  \text{LOF}_k(p) = \frac{1}{|N_k(p)|}
  \sum_{o \in N_k(p)} \frac{\text{lrd}_k(o)}{\text{lrd}_k(p)} .
  \label{eq:lof}
\end{equation}
A value of $\text{LOF}_k(p) \approx 1$ indicates that $p$ has a density
comparable to its neighbours (an inlier), whereas $\text{LOF}_k(p) \gg 1$
marks $p$ as an outlier sitting in a comparatively sparse region. In the
thesis, LOF is
configured with \texttt{n\_neighbors = 20} and the novelty-detection mode
(\texttt{novelty = True}) so that it can predict on unseen data. It is the
best-performing unsupervised model, with a precision of $88.28\%$ and a
recall of $95.13\%$, and is therefore chosen as the default third tier of
the system.

\subsection{Isolation Forest}

The Isolation Forest, proposed by Liu, Ting, and Zhou in
2008~\cite{liu2008}, approaches the problem from a different angle: instead
of profiling the normal points, the algorithm directly \emph{isolates} the
anomalous points. The idea rests on two properties of outliers: they are
the minority and they have attribute values markedly different from the
bulk of the data. By recursively and randomly partitioning the feature
space, anomalous points tend to be isolated after fewer splits,
corresponding to a smaller average depth in the tree.

Quantitatively, let $h(x)$ be the path length of a point $x$ (the number of
edges from the root to the leaf isolating it) and let $E(h(x))$ be its
average over all trees. The expected path length of an unsuccessful search in
a binary tree of $n$ points is used for normalisation,
\begin{equation}
  c(n) = 2H(n-1) - \frac{2(n-1)}{n},
  \qquad H(i) \approx \ln(i) + 0.5772,
  \label{eq:cn}
\end{equation}
where $0.5772$ is the Euler--Mascheroni constant. The anomaly score is then
\begin{equation}
  s(x,n) = 2^{-\,E(h(x))/c(n)} ,
  \label{eq:ifscore}
\end{equation}
so that $s \to 1$ when the average path length is very short (a likely
anomaly), and $s \approx 0.5$ when it is comparable to that of a normal
point. The advantage of the
Isolation Forest is its high speed and low memory demand thanks to its use
of sampling. In the thesis, this model attains high recall ($93.14\%$) but
a lower precision than LOF ($63.02\%$), so it serves as a comparison model
rather than the operational choice.

\subsection{One-Class SVM}

The One-Class SVM, proposed by Sch\"{o}lkopf et al.~\cite{scholkopf1999},
extends the support vector machine to the novelty-detection problem. The
algorithm learns a boundary enclosing the region of normal data in the
feature space; points falling outside this boundary are considered
anomalous. Mapping the inputs into a feature space through a function
$\phi(\cdot)$, it solves the optimisation problem
\begin{equation}
  \min_{\mathbf{w},\,\xi,\,\rho}\;
  \frac{1}{2}\lVert \mathbf{w}\rVert^{2}
  + \frac{1}{\nu n}\sum_{i=1}^{n}\xi_i - \rho
  \quad\text{s.t.}\quad
  \langle \mathbf{w}, \phi(x_i)\rangle \geq \rho - \xi_i,\;\; \xi_i \geq 0,
  \label{eq:ocsvm}
\end{equation}
and classifies a new point by the sign of the decision function
$f(x) = \operatorname{sign}\!\big(\langle \mathbf{w}, \phi(x)\rangle - \rho\big)$,
a negative sign indicating an anomaly. The parameter $\nu \in (0,1]$ upper-bounds
the fraction of training points allowed to fall outside the boundary. By
using a kernel function---typically the radial basis function (RBF) kernel
$K(x,x') = \exp(-\gamma\lVert x - x'\rVert^{2})$ as in the thesis---the
One-Class SVM can describe complex non-linear boundaries. Because the
algorithm is sensitive to feature scale, the thesis applies a scaling step
before training. The model attains a
recall of $91.38\%$ but a precision of only $62.27\%$, similar to the
Isolation Forest, reinforcing the conclusion that among the three one-class
algorithms, LOF is the most suitable for the HTTP log data of this problem.

% =====================================================================
\section{Feature Extraction and Evaluation Metrics}
\label{sec:features}

\subsection{Feature extraction from HTTP requests}

Because machine learning algorithms require numeric input, an essential
step is to transform each HTTP request into a feature vector. The thesis
does not use a bulky bag-of-words representation as in~\cite{pham2016} but
instead designs a compact $22$-dimensional vector that is shared by both
the supervised and the unsupervised tiers. Conceptually, these features
fall into three groups: structural features (length, path depth, parameter
count, and similar size measures), randomness-and-encoding features (the
Shannon entropy of the decoded payload and the number of encoding layers,
which expose obfuscation and evasion), and semantic-and-contextual features
(the signature score, the risk-keyword density, the suspicion of the
\texttt{User-Agent} and \texttt{Referer} headers, and the per-endpoint
unknown-parameter ratio). Using a shared feature space ensures consistency
and enables a fair comparison between the approaches, while keeping the
model compact and largely independent of the vocabulary of any specific
application. The precise composition of the vector and the extraction
procedure are presented in the next chapter.

\subsection{Evaluation metrics}

The performance of a detection system is usually assessed through the
confusion matrix, which comprises four quantities: the number of true
positives (TP, attacks detected correctly), true negatives (TN, clean
traffic recognised correctly), false positives (FP, clean traffic wrongly
flagged as an attack), and false negatives (FN, attacks that are missed).
From these four quantities, the thesis uses the following standard metrics,
defined consistently with Althubiti et al.~\cite{althubiti2017}.

\emph{Precision} is the proportion of attack predictions that are actually
correct, reflecting how ``clean'' the alerts are:
\begin{equation}
  \text{Precision} = \frac{TP}{TP + FP}.
  \label{eq:precision}
\end{equation}

\emph{Recall} is the proportion of actual attacks that the system catches,
reflecting the degree of missed detections:
\begin{equation}
  \text{Recall} = \frac{TP}{TP + FN}.
  \label{eq:recall}
\end{equation}

The \emph{F1 score} is the harmonic mean of precision and recall, providing
a single number that balances the two aspects:
\begin{equation}
  F_1 = \frac{2 \times \text{Precision} \times \text{Recall}}
             {\text{Precision} + \text{Recall}}.
  \label{eq:f1}
\end{equation}

In addition, the thesis emphasises the \emph{F2 score} as a metric
particularly suited to the intrusion detection problem:
\begin{equation}
  F_2 = \frac{5 \times \text{Precision} \times \text{Recall}}
             {4 \times \text{Precision} + \text{Recall}}.
  \label{eq:f2}
\end{equation}

The F2 score is a variant of the F-score family in which recall is weighted
twice as heavily as precision. This choice stems from the security context:
missing an attack (a false negative) may lead to a serious data breach,
whereas a false alarm (a false positive) costs the analyst only a few
minutes of investigation. Therefore, in the IDS context, prioritising
recall is reasonable, and the F2 score reflects that priority better than
the F1 score. The \texttt{\_safe\_metrics} function in
\texttt{apache\_log.py} computes all of the above metrics simultaneously for
every model configuration, forming the basis for the comparison table and
the evaluation charts in the subsequent chapters.

% =====================================================================
\section{Chapter Summary}
\label{sec:ch2summary}

This chapter has established the full foundation for moving on to the
design and implementation of the system. Section~\ref{sec:scope} delimited
the problem: detecting application-layer web attacks from Apache logs, with
a concrete catalogue of in-scope attack classes and a set of deliberately
excluded limitations. Section~\ref{sec:related} surveyed and critically
analysed a range of representative studies on the CSIC 2010 dataset,
showing that most research pursues a single approach and rarely combines
signature-based detection with machine learning; this gap is precisely the
motivation for the three-tier hybrid architecture of the thesis.
Sections~\ref{sec:ids} through~\ref{sec:features} presented the
foundational knowledge tightly bound to the system: the notion of an IDS
and its two detection paradigms, web attack types according to OWASP, the
regular-expression-based signature technique together with its
anti-evasion decoding mechanism, supervised learning algorithms (Random
Forest, Logistic Regression) and unsupervised one-class learning algorithms
(LOF, Isolation Forest, One-Class SVM), as well as the shared $22$-feature
extraction and the evaluation metrics centred on the F2 score. On the basis
of this knowledge and analysis, the next chapter will present the detailed
design of the three-tier architecture together with the smart consensus
decision mechanism, turning the theoretical principles laid out here into a
complete detection system.

% =====================================================================
% NOTE: The chapter-local bibliography has been removed. All references
% (pham2016, althubiti2017, nguyen2011, csic2010, rawashdeh2023, kim2020,
% vartouni2018, betarte2018, tsai2009, owasp2021, breiman2001,
% scholkopf1999, breunig2000, liu2008) now live in the global reference.bib
% and are printed once in the REFERENCE section via \printbibliography.
% =====================================================================

% --------------------- STANDALONE WRAPPER (begin) --------------------
\end{document}
% ---------------------- STANDALONE WRAPPER (end) ---------------------
